\section{Related Work}
\label{sec:related}

\subsection{Model Merging}
Model merging aggregates multiple task-specific expert neural networks sharing a pre-trained base into a unified multi-task model without joint retraining. Basic approaches like Task Arithmetic \cite{Ilharco2022TaskArithmetic} and Model Soups \cite{Wortsman2022Soups} merge task vectors linearly. To address parameter interference and sign conflicts, subsequent techniques like Git Re-Basin \cite{Ainsworth2022GitReBasin}, TIES-Merging \cite{Yadav2023TIES}, and DARE \cite{Yu2023DARE} align weights, resolve sign disagreements, or sparsify parameters.

To optimize merging parameters dynamically, AdaMerging \cite{Yang2023AdaMerging} introduced test-time adaptation of layer-wise coefficient weights using unlabeled data. Despite these advancements, existing model merging strategies primarily operate under the assumption of high-precision floating-point (FP32) weight spaces, leaving their performance under aggressive model compression unaddressed.

\subsection{Sharpness-Aware Minimization (SAM)}
Loss landscape geometry is a key determinant of generalization. SAM \cite{Foret2020SAM} biases optimization toward wider, flatter local minima by minimizing the maximum loss within an $l_2$ parameter neighborhood. Extensions like ASAM \cite{Kwon2021ASAM} and GSAM \cite{Zhuang2022GSAM} refine this with scale-invariance or surrogate gaps, while others adapt it for Vision Transformers (ViTs) \cite{Chen2021SAMViT}, which are prone to sharp minima.

While SAM's role in improving generalization is well-documented, its impact on downstream weight-space modifications is less understood. Recently, some studies have explored how SAM-trained experts interact when merged in full-precision. However, there is a total lack of research investigating how SAM-induced flatness influences the models' resilience to the severe non-linear perturbations introduced by post-training quantization.

\subsection{Post-Training Quantization (PTQ)}
Quantization compresses networks by mapping continuous weights to low-precision discrete representations, reducing memory and latency \cite{Jacob2018Quant}. PTQ is preferred because it avoids full retraining. For Transformers, standard PTQ incurs performance drops due to outlier channels, which methods like AdaRound \cite{Nagel2020AdaRound}, AWQ \cite{Lin2023AWQ}, and SmoothQuant \cite{Xiao2023SmoothQuant} mitigate by optimizing scale factors or rounding decisions on calibration subsets.

In model merging, quantization-aware merging (Q-Merge) \cite{Sterling2026QMerge} optimizes coefficients under quantization constraints using the Straight-Through Estimator (STE) \cite{Hubara2017Quantized}. Nevertheless, existing methods view the underlying experts as static inputs and fail to leverage pre-merging optimization to enhance robustness. FlatQ-Merge bridges this gap by directly manipulating expert flatness via SAM to demonstrate its deep empirical synergy with downstream PTQ and test-time coefficient adaptation.
